% Tripolar EEG — brief QC summary: TU2 & LS2 (Long felt-disc)
% Compile: pdflatex tu2_ls2_brief_report.tex (run twice if using cross-refs)
\documentclass[11pt]{article}
\usepackage[margin=1in]{geometry}
\usepackage{booktabs}
\usepackage{siunitx}
\usepackage{parskip}
\usepackage{xurl}
\usepackage{hyperref}

\sisetup{detect-all}

\title{Brief Analysis Report: Two Tripolar EEG Recordings\\{\large TU2 \& LS2 (Long felt-disc)}}
\author{Tripolar EEG study}
\date{\today}

\begin{document}
\maketitle

\section*{Purpose}
This memo reports pipeline outputs and automated QC metrics for two continuous recordings processed through the project single-subject workflow. Both \textbf{TU2} and \textbf{LS2} were collected with the \textbf{Long felt-disc} variant (same protocol class as HS and RK). Figures are in the Google Drive folder shared with Dr.\ Besio (see Output locations).

\section*{Artifacts generated (both participants)}
For each participant, the following figure set was produced: raw traces; ADC clipping summary; per-channel spectrograms (Ch1--Ch11); band-limited decomposition; alpha envelopes; alpha reactivity (eyes open vs.\ closed); open vs.\ closed PSD; PSD (alpha-focused band); correlation of alpha envelopes with the disc channel (Ch11); VEP comparison from the averaged file; five SSIM panels (one per TCRE pair) plus an SSIM summary bar chart; and a four-panel summary dashboard.

\section*{TU2 \textnormal{\textit{(Long felt-disc)}} }
\textit{Source files:} \texttt{TU2\_long\_felt\_TCRE\_4-7-2026} (\texttt{.eeg}, \texttt{.vhdr}, \texttt{.vmrk}; averaged VEP: \texttt{n\,=\,76} segments).

\begin{table}[h]
  \centering
  \small
  \begin{tabular}{@{}p{0.34\linewidth}p{0.60\linewidth}@{}}
    \toprule
    \textbf{Topic} & \textbf{Summary} \\
    \midrule
    Recording & ${\sim}6.2$\,min; 7 open/close epochs; 3 stimulation blocks; averaged VEP present ($n=76$). \\
    QC (automated) & Grade \textbf{MIXED}: ADC clipping mild, worst on Felt tEEG Ch5 (${\sim}0.99\%$ of samples; longest run ${\sim}0.65$\,s). Near-flat windows minimal (${\sim}0.27\%$ on Ch5). \\
    Alpha SNR & Median ${\sim}2.6$\,dB; best channel ${\sim}6.6$\,dB. \\
    Alpha reactivity (closed/open $\alpha$ power) & Median ${\sim}2.63\times$ (the strongest median in the long felt-disc cohort); paste/disc: Ch9 ${\sim}6.9\times$, Ch10 ${\sim}8.2\times$, Ch11 ${\sim}8.0\times$. \\
    $\alpha$ envelope vs.\ disc (Ch11) & Paste Ch9--Ch10: $r \approx 0.995$--$0.999$ (exceptionally high). Median across all channels ${\sim}0.51$. \\
    Spectrogram SSIM (tEEG vs.\ eEEG, same TCRE) & Felt pairs \#1--\#4: 0.27--0.34 (low agreement). Paste pair \#5: ${\sim}0.66$ (moderate--high). \\
    Peak $\alpha$ (eyes closed, Welch) & Paste tEEG Ch9 ${\sim}11.0$\,Hz; felt tEEG Ch1 ${\sim}8.3$\,Hz (channel-dependent). \\
    VEP peak-to-peak (averaged waveform) & Felt tEEG Ch1 ${\sim}243$\,$\mu$V; paste Ch9 ${\sim}101$\,$\mu$V; Ch10--Ch11 ${\sim}7$\,$\mu$V. \\
    \bottomrule
  \end{tabular}
\end{table}

\noindent\textit{Interpretation (brief):} TU2 is the \textbf{cleanest recording} in the long felt-disc group by QC grade (MIXED vs.\ POOR for HS, RK, and LS2). Clipping is minimal; paste disc correlations approach unity ($r > 0.99$); and alpha reactivity is strong on both paste and several felt channels. Felt spectrogram SSIM remains low ($0.27$--$0.34$), consistent with the rest of the cohort.

\section*{LS2 \textnormal{\textit{(Long felt-disc)}}}
\textit{Source files:} \texttt{LS2-long-felt-TCRE-4-7-2026} (\texttt{.eeg}, \texttt{.vhdr}, \texttt{.vmrk}; averaged VEP: \texttt{n\,=\,74} segments).

\begin{table}[h]
  \centering
  \small
  \begin{tabular}{@{}p{0.34\linewidth}p{0.60\linewidth}@{}}
    \toprule
    \textbf{Topic} & \textbf{Summary} \\
    \midrule
    Recording & ${\sim}6.3$\,min; 7 open/close epochs; 3 stimulation blocks; averaged VEP ($n=74$). \\
    QC (automated) & Grade POOR: ADC clipping worst on Felt tEEG Ch1 (${\sim}21.9\%$ of samples; longest run ${\sim}0.48$\,s). Near-flat fraction minimal (${\sim}0.27\%$ on Ch1). \\
    Alpha SNR & Median ${\sim}2.9$\,dB; best channel ${\sim}7.0$\,dB (best in the long felt-disc group). \\
    Alpha reactivity & Median ${\sim}1.81\times$; paste/disc strong: Ch9 ${\sim}4.3\times$, Ch10 ${\sim}4.7\times$, Ch11 ${\sim}5.5\times$. \\
    $\alpha$ envelope vs.\ disc & Paste Ch9--10: $r \approx 0.977$--$0.983$. Median across all channels ${\sim}0.59$ (highest in group). \\
    Spectrogram SSIM & Felt \#1--\#4: 0.28--0.34; paste \#5: ${\sim}0.79$ (highest paste SSIM in the long felt-disc cohort). \\
    Peak $\alpha$ (eyes closed) & Paste Ch9 ${\sim}10.7$\,Hz; felt Ch1 ${\sim}11.5$\,Hz. \\
    VEP peak-to-peak & Felt tEEG Ch1 ${\sim}329$\,$\mu$V; paste Ch9 ${\sim}106$\,$\mu$V; Ch10--11 ${\sim}10$--$14$\,$\mu$V. \\
    \bottomrule
  \end{tabular}
\end{table}

\noindent\textit{Interpretation (brief):} LS2 shows the \textbf{best spectral quality} (highest alpha SNR and best paste SSIM) in the long felt-disc group despite moderate clipping on Felt tEEG Ch1 (${\sim}22\%$). Paste/disc metrics are very strong; felt SSIM is still low ($0.28$--$0.34$), again consistent with the group.

\section*{Side-by-side comparison}
Both participants share the Long felt-disc protocol. \textbf{TU2} has negligible clipping ($<1\%$) and the strongest alpha reactivity (median ${\sim}2.6\times$), making it the best-quality recording in the long felt-disc group overall. \textbf{LS2} carries heavier Ch1 clipping (${\sim}22\%$) but compensates with the highest paste SSIM (${\sim}0.79$) and alpha SNR (${\sim}2.9$\,dB median) in the group. Across both, the Felt tEEG vs.\ eEEG spectrogram agreement (SSIM $0.27$--$0.34$) is consistently lower than the paste pair ($0.66$--$0.79$), echoing the same pattern seen in HS and RK.

\section*{Output locations}
\noindent\textbf{Shared folder (Google Drive, Dr.\ Besio):}\\
\href{https://drive.google.com/drive/folders/1_dWUq1SyyP8dv59GeOEPtZUoBWs0N9XL?usp=drive_link}{drive.google.com/drive/folders/1\_dWUq1SyyP8dv59GeOEPtZUoBWs0N9XL}

\end{document}
